\documentclass[12pt,a4paper]{article}

% Paquetes
\usepackage[utf8]{inputenc}
\usepackage[spanish]{babel}
\usepackage{geometry}
\geometry{margin=2.5cm}
\usepackage{hyperref}
\usepackage{xcolor}
\usepackage{graphicx}
\usepackage{fancyhdr}
\usepackage{titlesec}
\usepackage{enumitem}
\usepackage{booktabs}
\usepackage{array}

% Colores personalizados
\definecolor{azulgit}{RGB}{24,80,50}
\definecolor{grisclaro}{RGB}{245,245,245}

% Configuración de hipervínculos
\hypersetup{
    colorlinks=true,
    linkcolor=azulgit,
    urlcolor=azulgit,
    citecolor=azulgit
}

% Encabezado y pie de página
\pagestyle{fancy}
\fancyhf{}
\fancyhead[L]{\textcolor{azulgit}{\small Programadores Por La Paz}}
\fancyhead[R]{\textcolor{azulgit}{\small Semana 3}}
\fancyfoot[C]{\thepage}
\renewcommand{\headrulewidth}{0.4pt}

% Títulos
\titleformat{\section}
  {\normalfont\Large\bfseries\color{azulgit}}
  {\thesection}{1em}{}
\titleformat{\subsection}
  {\normalfont\large\bfseries\color{azulgit!80!black}}
  {\thesubsection}{1em}{}

\begin{document}

% Título
\begin{titlepage}
    \centering
    \vspace*{2cm}
    
    {\Huge\bfseries\textcolor{azulgit}{DOCUMENTO DE PRESENTACIÓN}\\[0.5cm]}
    {\Large\textcolor{azulgit!70!black}{Tarea Semana 3}}\\[1cm]
    
    \vspace{1cm}
    
    {\Large Git y Control de Versiones}\\[2cm]
    
    \vspace{2cm}
    
    {\large\textbf{Programa:}}\\
    {\large Programadores Por La Paz}\\[1.5cm]
    
    {\large\textbf{Universidad de Cartagena}}\\[2cm]
    
    \vfill
    
    {\large\textbf{Estudiante:}}\\
    {\Large Devin Polanco Romero}\\[1cm]
    
    {\large Junio 2026}
    
\end{titlepage}

% Índice
\tableofcontents
\newpage

% Contenido principal
\section{Información General}

\begin{table}[h]
\centering
\begin{tabular}{|>{\bfseries}l|l|}
\hline
\textbf{Campo} & \textbf{Detalle} \\
\hline
Estudiante & Devin Polanco Romero \\
\hline
Programa & Programadores Por La Paz \\
\hline
Institución & Universidad de Cartagena \\
\hline
Actividad & Tarea Semana 3 \\
\hline
Tema & Git y Control de Versiones \\
\hline
Fecha & Junio 2026 \\
\hline
\end{tabular}
\caption{Información de la actividad}
\end{table}

\section{Enlace del Repositorio}

El desarrollo de la actividad fue publicado en el siguiente repositorio de GitHub:

\vspace{0.5cm}

\begin{center}
\large
\url{https://github.com/DevCorpSoftwareFactory/Programadores_Para_La_Paz_Semana_3}
\end{center}

\vspace{0.5cm}

\noindent La carpeta \texttt{semana3} contiene todos los archivos solicitados para la entrega de la actividad.

\section{Objetivo de la Semana}

Comprender los conceptos básicos de \textbf{Git} y \textbf{control de versiones}, aprender a registrar cambios en un proyecto, entender la importancia de la \textbf{trazabilidad} en el desarrollo de software y continuar utilizando el repositorio como espacio de trabajo y entrega.

\section{Desarrollo de la Actividad}

\subsection{Preguntas de Selección Múltiple}

Se respondieron cuatro preguntas sobre conceptos fundamentales de Git:

\begin{enumerate}[label=\textbf{\arabic*.}]
    \item \textbf{¿Qué es Git?} \\
    Respuesta: Un sistema de control de versiones para registrar cambios en proyectos.
    
    \item \textbf{¿Cuál es la función principal de un repositorio?} \\
    Respuesta: Registrar y organizar los archivos y cambios de un proyecto.
    
    \item \textbf{¿Qué hace el comando \texttt{git add .}?} \\
    Respuesta: Agrega archivos al área de preparación para registrar cambios.
    
    \item \textbf{¿Qué hace el comando \texttt{git commit}?} \\
    Respuesta: Guarda un registro de cambios realizados en el proyecto.
\end{enumerate}

\subsection{Conceptos Básicos de Git}

Se explicaron los siguientes conceptos fundamentales:

\begin{description}
    \item[Repositorio] Espacio donde se guardan los archivos y todos los cambios realizados en un proyecto. Puede ser local (en la computador) o remoto (en la nube como GitHub).
    
    \item[Commit] Registro de cambios realizado en un momento específico del proyecto. Cada commit tiene un identificador único y permite volver a versiones anteriores.
    
    \item[Push] Acción de enviar los cambios desde la máquina local hacia GitHub, permitiendo la colaboración y el trabajo en equipo.
\end{description}

\subsection{Reflexión}

Se desarrolló una reflexión de 8 a 10 líneas sobre la importancia de mantener un registro de cambios en proyectos de software usando herramientas como Git, destacando:

\begin{itemize}
    \item La necesidad de tener un historial completo y organizado
    \item La posibilidad de volver a versiones anteriores
    \item La facilitación del trabajo colaborativo en equipo
    \item La importancia de la trazabilidad en el desarrollo de software
\end{itemize}

\section{Archivos Entregados}

La carpeta \texttt{semana3} contiene los siguientes archivos:

\begin{table}[h]
\centering
\begin{tabular}{|l|l|}
\hline
\textbf{Archivo} & \textbf{Contenido} \\
\hline
preguntas-semana3.txt & Respuestas de selección múltiple \\
\hline
conceptos-git.txt & Explicación de repositorio, commit y push \\
\hline
reflexion-semana3.txt & Reflexión sobre la importancia de Git \\
\hline
\end{tabular}
\caption{Archivos de la actividad}
\end{table}

\section{Conclusiones}

\begin{itemize}
    \item Se comprenderon los conceptos básicos de Git y su importancia en el desarrollo de software.
    \item Se practicó el uso de comandos fundamentales como \texttt{git add}, \texttt{git commit} y \texttt{git push}.
    \item Se reconoció la trazabilidad como un elemento esencial para la calidad y organización de proyectos de software.
    \item Se fortaleció el uso de GitHub como herramienta de publicación y seguimiento del avance técnico.
\end{itemize}

\section{Referencias}

\begin{itemize}
    \item Documentación oficial de Git: \url{https://git-scm.com/doc}
    \item Documentación de GitHub: \url{https://docs.github.com}
    \item Libro Pro Git: \url{https://git-scm.com/book/es/v2}
\end{itemize}

\end{document}
